
\documentclass[10pt,letterpaper]{article}
\usepackage[T1]{fontenc}
\usepackage[utf8]{inputenc}
\usepackage{lmodern}
\usepackage[margin=0.72in,headheight=14pt]{geometry}
\usepackage{xcolor}
\usepackage{hyperref}
\usepackage{bookmark}
\usepackage{enumitem}
\usepackage{booktabs}
\usepackage{longtable}
\usepackage{array}
\usepackage{ragged2e}
\usepackage{titlesec}
\usepackage{fancyhdr}
\usepackage{microtype}
\usepackage{listings}
\usepackage{needspace}
\usepackage{tocloft}
\usepackage{upquote}
\definecolor{HouseBlue}{HTML}{1F4E79}
\definecolor{HouseGray}{HTML}{4F5B62}
\definecolor{HouseLight}{HTML}{EAF2F8}
\definecolor{CodeBack}{HTML}{F7F8FA}
\definecolor{CodeFrame}{HTML}{C9D1D9}
\definecolor{CodeKeyword}{HTML}{1F4E79}
\definecolor{CodeComment}{HTML}{5F6F7A}
\definecolor{CodeString}{HTML}{8A3B12}
\hypersetup{colorlinks=true,linkcolor=HouseBlue,urlcolor=HouseBlue,citecolor=HouseBlue,pdfauthor={ChatGPT},pdftitle={AI for Games: Algorithms as C-like Pseudocode}}
\pagestyle{fancy}
\fancyhf{}
\lhead{\small AI for Games --- Algorithms as C-like Pseudocode}
\rhead{\small\thepage}
\renewcommand{\headrulewidth}{0.4pt}
\setlength{\parindent}{0pt}
\setlength{\parskip}{5pt plus 1pt}
\setlist[itemize]{leftmargin=1.25em,itemsep=2pt,topsep=2pt}
\setlist[enumerate]{leftmargin=1.5em,itemsep=2pt,topsep=2pt}
\titleformat{\section}{\Large\bfseries\color{HouseBlue}}{}{0pt}{}
\titleformat{\subsection}{\large\bfseries\color{HouseGray}}{}{0pt}{}
\titleformat{\subsubsection}{\normalsize\bfseries\color{HouseGray}}{}{0pt}{}
\titlespacing*{\section}{0pt}{2.2ex plus .4ex minus .2ex}{1.0ex}
\titlespacing*{\subsection}{0pt}{1.8ex plus .3ex minus .2ex}{0.7ex}
\newcolumntype{L}[1]{>{\RaggedRight\arraybackslash}p{#1}}
\lstdefinelanguage{PseudoC}{
  morekeywords={typedef,enum,struct,union,return,if,else,for,while,do,case,switch,break,continue,bool,true,false,NULL,void,int,long,char,float,double,string,const,static,foreach,in,new,delete,insert,remove,push,pop,emit,try,catch,throw,private,public,Map,List,Set,Queue,Stack,PriorityQueue,Vec2,Vec3,Vector,Matrix,Node,Edge,Graph,Path,State,Action,Goal,Behavior,Task,Status,Success,Failure,Running,GameState,Character,Decision,Rule,Condition,Event,Sensor,Region,Value},
  sensitive=true,
  morecomment=[l]{//},
  morestring=[b]",
}
\lstset{
  language=PseudoC,
  basicstyle=\ttfamily\scriptsize,
  keywordstyle=\bfseries\color{CodeKeyword},
  commentstyle=\itshape\color{CodeComment},
  stringstyle=\color{CodeString},
  backgroundcolor=\color{CodeBack},
  frame=single,
  rulecolor=\color{CodeFrame},
  framerule=0.35pt,
  breaklines=true,
  breakatwhitespace=false,
  columns=fullflexible,
  keepspaces=true,
  showstringspaces=false,
  tabsize=4,
  xleftmargin=0.6em,
  xrightmargin=0.3em,
  aboveskip=0.7em,
  belowskip=0.7em,
}
\newcommand{\handouttitle}{AI for Games: Algorithms as C-like Pseudocode}
\sloppy
\begin{document}
\pagenumbering{gobble}
\begin{titlepage}
\centering
\vspace*{0.85in}
{\Huge\bfseries\color{HouseBlue}\handouttitle\par}
\vspace{0.22in}
{\Large Systematic Algorithm Extraction and Reconstruction\par}
\vspace{0.34in}
{\large Source analyzed: Ian Millington, \emph{AI for Games}, Third Edition, CRC Press, 2019\par}
\vfill
\begin{minipage}{0.88\textwidth}
\small
This handout rewrites the algorithmic content from the uploaded game-AI text as independent C-like pseudocode. It is organized as an implementation-oriented reference for movement, steering behaviors, pathfinding, decision making, tactical and strategic AI, learning, procedural content generation, board-game search, execution management, event handling, polling, and sense management. The pseudocode is normalized for study, code review, engine design, and implementation planning; it is not a source-code clone.
\end{minipage}
\vfill
{\large Generated \today\par}
\end{titlepage}
\clearpage
\pagenumbering{roman}
\tableofcontents
\clearpage
\pagenumbering{arabic}


\textbf{Source analyzed:} Ian Millington, \emph{AI for Games}, Third Edition, CRC Press, 2019.
\textbf{Output purpose:} compact algorithm-study reference for implementation planning, code review, and engine design.
\textbf{Method:} algorithms, pseudo-code sections, and code-like data structures were systematically reduced to independent C-like pseudocode. This is not a transcription of the book's source listings; names and control flow are normalized for C/C++-style implementation.

\medskip\hrule\medskip

\Needspace{5\baselineskip}
\subsection*{0. Shared Data Structures and Utility Routines}
\addcontentsline{toc}{subsection}{0. Shared Data Structures and Utility Routines}

\begin{lstlisting}
typedef struct Vec3 {
    float x, y, z;
} Vec3;

typedef struct Kinematic {
    Vec3 position;
    float orientation;      // yaw in radians for 2.5D movement
    Vec3 velocity;
    float rotation;         // angular velocity
} Kinematic;

typedef struct SteeringOutput {
    Vec3 linear;
    float angular;
    bool valid;
} SteeringOutput;

typedef struct KinematicOutput {
    Vec3 velocity;
    float rotation;
    bool valid;
} KinematicOutput;

float length(Vec3 v);
Vec3 normalize(Vec3 v);
Vec3 add(Vec3 a, Vec3 b);
Vec3 sub(Vec3 a, Vec3 b);
Vec3 scale(Vec3 v, float s);
float dot(Vec3 a, Vec3 b);
float random01(void);           // [0, 1)
float randomSigned(void);       // [-1, 1]
float clamp(float x, float lo, float hi);
float wrapPi(float radians);    // normalize angle to [-pi, pi]
\end{lstlisting}

\begin{lstlisting}
float randomBinomial(void) {
    return random01() - random01();
}

Vec3 orientationAsVector(float orientation) {
    Vec3 v = { sinf(orientation), 0.0f, cosf(orientation) };
    return v;
}

float newOrientation(float current, Vec3 velocity) {
    if (length(velocity) <= 0.0f) return current;
    return atan2f(velocity.x, velocity.z);
}
\end{lstlisting}

\begin{lstlisting}
void updateKinematic(Kinematic *k, SteeringOutput s, float dt) {
    if (!s.valid) return;
    k->position = add(k->position, scale(k->velocity, dt));
    k->orientation += k->rotation * dt;
    k->velocity = add(k->velocity, scale(s.linear, dt));
    k->rotation += s.angular * dt;
}
\end{lstlisting}

\medskip\hrule\medskip

\clearpage
\section*{Part I -- Movement Algorithms}
\addcontentsline{toc}{section}{Part I -- Movement Algorithms}

\Needspace{5\baselineskip}
\subsection*{1. Kinematic Seek}
\addcontentsline{toc}{subsection}{1. Kinematic Seek}

\begin{lstlisting}
KinematicOutput kinematicSeek(Kinematic character, Kinematic target, float maxSpeed) {
    KinematicOutput out = {0};
    Vec3 direction = sub(target.position, character.position);
    if (length(direction) <= 0.0f) return out;
    out.velocity = scale(normalize(direction), maxSpeed);
    out.rotation = 0.0f;
    out.valid = true;
    return out;
}
\end{lstlisting}

\Needspace{5\baselineskip}
\subsection*{2. Kinematic Flee}
\addcontentsline{toc}{subsection}{2. Kinematic Flee}

\begin{lstlisting}
KinematicOutput kinematicFlee(Kinematic character, Kinematic target, float maxSpeed) {
    KinematicOutput out = {0};
    Vec3 direction = sub(character.position, target.position);
    if (length(direction) <= 0.0f) return out;
    out.velocity = scale(normalize(direction), maxSpeed);
    out.rotation = 0.0f;
    out.valid = true;
    return out;
}
\end{lstlisting}

\Needspace{5\baselineskip}
\subsection*{3. Kinematic Arrive}
\addcontentsline{toc}{subsection}{3. Kinematic Arrive}

\begin{lstlisting}
KinematicOutput kinematicArrive(
    Kinematic character,
    Kinematic target,
    float maxSpeed,
    float radius,
    float timeToTarget
) {
    KinematicOutput out = {0};
    Vec3 direction = sub(target.position, character.position);
    float distance = length(direction);
    if (distance < radius) return out;

    Vec3 desired = scale(direction, 1.0f / timeToTarget);
    if (length(desired) > maxSpeed) desired = scale(normalize(desired), maxSpeed);

    out.velocity = desired;
    out.rotation = 0.0f;
    out.valid = true;
    return out;
}
\end{lstlisting}

\Needspace{5\baselineskip}
\subsection*{4. Kinematic Wander}
\addcontentsline{toc}{subsection}{4. Kinematic Wander}

\begin{lstlisting}
KinematicOutput kinematicWander(Kinematic character, float maxSpeed, float maxRotation) {
    KinematicOutput out = {0};
    out.velocity = scale(orientationAsVector(character.orientation), maxSpeed);
    out.rotation = randomBinomial() * maxRotation;
    out.valid = true;
    return out;
}
\end{lstlisting}

\Needspace{5\baselineskip}
\subsection*{5. Dynamic Seek}
\addcontentsline{toc}{subsection}{5. Dynamic Seek}

\begin{lstlisting}
SteeringOutput seek(Kinematic character, Kinematic target, float maxAccel) {
    SteeringOutput out = {0};
    Vec3 direction = sub(target.position, character.position);
    if (length(direction) <= 0.0f) return out;
    out.linear = scale(normalize(direction), maxAccel);
    out.angular = 0.0f;
    out.valid = true;
    return out;
}
\end{lstlisting}

\Needspace{5\baselineskip}
\subsection*{6. Dynamic Flee}
\addcontentsline{toc}{subsection}{6. Dynamic Flee}

\begin{lstlisting}
SteeringOutput flee(Kinematic character, Kinematic target, float maxAccel) {
    SteeringOutput out = {0};
    Vec3 direction = sub(character.position, target.position);
    if (length(direction) <= 0.0f) return out;
    out.linear = scale(normalize(direction), maxAccel);
    out.angular = 0.0f;
    out.valid = true;
    return out;
}
\end{lstlisting}

\Needspace{5\baselineskip}
\subsection*{7. Arrive Steering}
\addcontentsline{toc}{subsection}{7. Arrive Steering}

\begin{lstlisting}
SteeringOutput arrive(
    Kinematic character,
    Kinematic target,
    float maxSpeed,
    float maxAccel,
    float targetRadius,
    float slowRadius,
    float timeToTarget
) {
    SteeringOutput out = {0};
    Vec3 direction = sub(target.position, character.position);
    float distance = length(direction);
    if (distance < targetRadius) return out;

    float targetSpeed = (distance > slowRadius)
        ? maxSpeed
        : maxSpeed * distance / slowRadius;

    Vec3 targetVelocity = scale(normalize(direction), targetSpeed);
    out.linear = scale(sub(targetVelocity, character.velocity), 1.0f / timeToTarget);
    if (length(out.linear) > maxAccel) out.linear = scale(normalize(out.linear), maxAccel);
    out.angular = 0.0f;
    out.valid = true;
    return out;
}
\end{lstlisting}

\Needspace{5\baselineskip}
\subsection*{8. Align Steering}
\addcontentsline{toc}{subsection}{8. Align Steering}

\begin{lstlisting}
SteeringOutput align(
    Kinematic character,
    Kinematic target,
    float maxRotation,
    float maxAngularAccel,
    float targetRadius,
    float slowRadius,
    float timeToTarget
) {
    SteeringOutput out = {0};
    float rotation = wrapPi(target.orientation - character.orientation);
    float rotationSize = fabsf(rotation);
    if (rotationSize < targetRadius) return out;

    float targetRotation = (rotationSize > slowRadius)
        ? maxRotation
        : maxRotation * rotationSize / slowRadius;
    targetRotation *= rotation / rotationSize;

    out.angular = (targetRotation - character.rotation) / timeToTarget;
    out.angular = clamp(out.angular, -maxAngularAccel, maxAngularAccel);
    out.linear = (Vec3){0,0,0};
    out.valid = true;
    return out;
}
\end{lstlisting}

\Needspace{5\baselineskip}
\subsection*{9. Velocity Matching}
\addcontentsline{toc}{subsection}{9. Velocity Matching}

\begin{lstlisting}
SteeringOutput velocityMatch(Kinematic character, Kinematic target, float maxAccel, float timeToTarget) {
    SteeringOutput out = {0};
    out.linear = scale(sub(target.velocity, character.velocity), 1.0f / timeToTarget);
    if (length(out.linear) > maxAccel) out.linear = scale(normalize(out.linear), maxAccel);
    out.angular = 0.0f;
    out.valid = true;
    return out;
}
\end{lstlisting}

\Needspace{5\baselineskip}
\subsection*{10. Pursue and Evade}
\addcontentsline{toc}{subsection}{10. Pursue and Evade}

\begin{lstlisting}
Kinematic predictTarget(Kinematic character, Kinematic target, float maxPrediction) {
    Vec3 direction = sub(target.position, character.position);
    float distance = length(direction);
    float speed = length(character.velocity);
    float prediction = (speed <= distance / maxPrediction) ? maxPrediction : distance / speed;
    target.position = add(target.position, scale(target.velocity, prediction));
    return target;
}

SteeringOutput pursue(Kinematic character, Kinematic target, float maxPrediction, float maxAccel) {
    Kinematic future = predictTarget(character, target, maxPrediction);
    return seek(character, future, maxAccel);
}

SteeringOutput evade(Kinematic character, Kinematic target, float maxPrediction, float maxAccel) {
    Kinematic future = predictTarget(character, target, maxPrediction);
    return flee(character, future, maxAccel);
}
\end{lstlisting}

\Needspace{5\baselineskip}
\subsection*{11. Face and Look-Where-You-Are-Going}
\addcontentsline{toc}{subsection}{11. Face and Look-Where-You-Are-Going}

\begin{lstlisting}
SteeringOutput face(
    Kinematic character,
    Kinematic target,
    float maxRotation,
    float maxAngularAccel,
    float targetRadius,
    float slowRadius,
    float timeToTarget
) {
    Vec3 direction = sub(target.position, character.position);
    SteeringOutput out = {0};
    if (length(direction) <= 0.0f) return out;

    Kinematic orientationTarget = target;
    orientationTarget.orientation = atan2f(direction.x, direction.z);
    return align(character, orientationTarget, maxRotation, maxAngularAccel,
                 targetRadius, slowRadius, timeToTarget);
}

SteeringOutput lookWhereGoing(
    Kinematic character,
    float maxRotation,
    float maxAngularAccel,
    float targetRadius,
    float slowRadius,
    float timeToTarget
) {
    Kinematic target = character;
    if (length(character.velocity) <= 0.0f) return (SteeringOutput){0};
    target.orientation = atan2f(character.velocity.x, character.velocity.z);
    return align(character, target, maxRotation, maxAngularAccel,
                 targetRadius, slowRadius, timeToTarget);
}
\end{lstlisting}

\Needspace{5\baselineskip}
\subsection*{12. Dynamic Wander}
\addcontentsline{toc}{subsection}{12. Dynamic Wander}

\begin{lstlisting}
typedef struct WanderState {
    float wanderOrientation;
} WanderState;

SteeringOutput wander(
    Kinematic character,
    WanderState *state,
    float wanderOffset,
    float wanderRadius,
    float wanderRate,
    float maxAccel
) {
    state->wanderOrientation += randomBinomial() * wanderRate;
    float targetOrientation = state->wanderOrientation + character.orientation;

    Kinematic target = {0};
    target.position = add(character.position, scale(orientationAsVector(character.orientation), wanderOffset));
    target.position = add(target.position, scale(orientationAsVector(targetOrientation), wanderRadius));
    return face(character, target, /*maxRotation*/9999, /*maxAngularAccel*/9999,
                /*targetRadius*/0.01f, /*slowRadius*/1.0f, /*timeToTarget*/0.1f);
}
\end{lstlisting}

\Needspace{5\baselineskip}
\subsection*{13. Path Following}
\addcontentsline{toc}{subsection}{13. Path Following}

\begin{lstlisting}
typedef struct Path {
    float (*getParam)(struct Path *p, Vec3 position, float lastParam);
    Vec3  (*getPosition)(struct Path *p, float param);
} Path;

SteeringOutput followPath(
    Kinematic character,
    Path *path,
    float *pathParam,
    float pathOffset,
    float maxAccel
) {
    *pathParam = path->getParam(path, character.position, *pathParam);
    float targetParam = *pathParam + pathOffset;
    Kinematic target = {0};
    target.position = path->getPosition(path, targetParam);
    return seek(character, target, maxAccel);
}
\end{lstlisting}

\Needspace{5\baselineskip}
\subsection*{14. Separation}
\addcontentsline{toc}{subsection}{14. Separation}

\begin{lstlisting}
SteeringOutput separation(
    Kinematic character,
    Kinematic *targets,
    int count,
    float threshold,
    float decayCoefficient,
    float maxAccel
) {
    SteeringOutput out = {0};
    for (int i = 0; i < count; ++i) {
        Vec3 direction = sub(character.position, targets[i].position);
        float distance = length(direction);
        if (distance <= 0.0f || distance >= threshold) continue;
        float strength = fminf(decayCoefficient / (distance * distance), maxAccel);
        out.linear = add(out.linear, scale(normalize(direction), strength));
    }
    out.angular = 0.0f;
    out.valid = length(out.linear) > 0.0f;
    return out;
}
\end{lstlisting}

\Needspace{5\baselineskip}
\subsection*{15. Collision Avoidance}
\addcontentsline{toc}{subsection}{15. Collision Avoidance}

\begin{lstlisting}
SteeringOutput collisionAvoidance(
    Kinematic character,
    Kinematic *targets,
    int count,
    float radius,
    float maxAccel
) {
    float shortestTime = INFINITY;
    int first = -1;
    Vec3 firstRelativePos = {0}, firstRelativeVel = {0};

    for (int i = 0; i < count; ++i) {
        Vec3 relPos = sub(targets[i].position, character.position);
        Vec3 relVel = sub(targets[i].velocity, character.velocity);
        float relSpeed2 = dot(relVel, relVel);
        if (relSpeed2 <= 0.0f) continue;

        float time = -dot(relPos, relVel) / relSpeed2;
        float distanceAtClosest = length(add(relPos, scale(relVel, time)));
        if (time > 0.0f && distanceAtClosest < 2.0f * radius && time < shortestTime) {
            shortestTime = time;
            first = i;
            firstRelativePos = relPos;
            firstRelativeVel = relVel;
        }
    }

    SteeringOutput out = {0};
    if (first < 0) return out;
    Vec3 avoidDirection = add(firstRelativePos, scale(firstRelativeVel, shortestTime));
    out.linear = scale(normalize(scale(avoidDirection, -1.0f)), maxAccel);
    out.valid = true;
    return out;
}
\end{lstlisting}

\Needspace{5\baselineskip}
\subsection*{16. Obstacle and Wall Avoidance}
\addcontentsline{toc}{subsection}{16. Obstacle and Wall Avoidance}

\begin{lstlisting}
typedef struct Collision {
    bool hit;
    Vec3 position;
    Vec3 normal;
} Collision;

typedef struct CollisionDetector {
    Collision (*raycast)(struct CollisionDetector *detector, Vec3 start, Vec3 end);
} CollisionDetector;

SteeringOutput obstacleAvoidance(
    Kinematic character,
    CollisionDetector *detector,
    float avoidDistance,
    float lookahead,
    float maxAccel
) {
    SteeringOutput out = {0};
    Vec3 ray = scale(normalize(character.velocity), lookahead);
    Collision c = detector->raycast(detector, character.position, add(character.position, ray));
    if (!c.hit) return out;

    Kinematic target = {0};
    target.position = add(c.position, scale(c.normal, avoidDistance));
    return seek(character, target, maxAccel);
}
\end{lstlisting}

\Needspace{5\baselineskip}
\subsection*{17. Weighted Blended Steering}
\addcontentsline{toc}{subsection}{17. Weighted Blended Steering}

\begin{lstlisting}
typedef SteeringOutput (*SteeringBehavior)(void *ctx);

typedef struct WeightedBehavior {
    SteeringBehavior behavior;
    void *ctx;
    float weight;
} WeightedBehavior;

SteeringOutput blendedSteering(WeightedBehavior *items, int count, float maxAccel, float maxAngularAccel) {
    SteeringOutput out = {0};
    for (int i = 0; i < count; ++i) {
        SteeringOutput s = items[i].behavior(items[i].ctx);
        if (!s.valid) continue;
        out.linear = add(out.linear, scale(s.linear, items[i].weight));
        out.angular += s.angular * items[i].weight;
        out.valid = true;
    }
    if (length(out.linear) > maxAccel) out.linear = scale(normalize(out.linear), maxAccel);
    out.angular = clamp(out.angular, -maxAngularAccel, maxAngularAccel);
    return out;
}
\end{lstlisting}

\Needspace{5\baselineskip}
\subsection*{18. Priority Steering}
\addcontentsline{toc}{subsection}{18. Priority Steering}

\begin{lstlisting}
SteeringOutput prioritySteering(WeightedBehavior *groups, int count, float epsilon) {
    for (int i = 0; i < count; ++i) {
        SteeringOutput s = groups[i].behavior(groups[i].ctx);
        if (s.valid && (length(s.linear) > epsilon || fabsf(s.angular) > epsilon)) return s;
    }
    return (SteeringOutput){0};
}
\end{lstlisting}

\Needspace{5\baselineskip}
\subsection*{19. Steering Pipeline}
\addcontentsline{toc}{subsection}{19. Steering Pipeline}

\begin{lstlisting}
typedef struct Goal { Vec3 position; float orientation; } Goal;
typedef bool (*ConstraintCheck)(Goal g, void *ctx);
typedef Goal (*ConstraintSuggest)(Goal g, void *ctx);
typedef Goal (*Targeter)(void *ctx);
typedef Goal (*Decomposer)(Goal g, void *ctx);
typedef SteeringOutput (*Actuator)(Goal g, void *ctx);

SteeringOutput steeringPipeline(
    Targeter targeter,
    Decomposer decomposer,
    ConstraintCheck *checks,
    ConstraintSuggest *suggestions,
    void **constraintCtx,
    int constraintCount,
    Actuator actuator,
    void *ctx
) {
    Goal goal = targeter(ctx);
    goal = decomposer(goal, ctx);

    for (int i = 0; i < constraintCount; ++i) {
        if (checks[i](goal, constraintCtx[i])) {
            goal = suggestions[i](goal, constraintCtx[i]);
            i = -1; // restart constraint validation after a repair
        }
    }
    return actuator(goal, ctx);
}
\end{lstlisting}

\Needspace{5\baselineskip}
\subsection*{20. Ballistic Firing Solution}
\addcontentsline{toc}{subsection}{20. Ballistic Firing Solution}

\begin{lstlisting}
bool firingSolution(Vec3 shooter, Vec3 target, float projectileSpeed, Vec3 gravity, Vec3 *velocityOut) {
    Vec3 delta = sub(target, shooter);
    float g = length(gravity);
    float xz = sqrtf(delta.x * delta.x + delta.z * delta.z);
    float y = delta.y;
    float v2 = projectileSpeed * projectileSpeed;
    float root = v2 * v2 - g * (g * xz * xz + 2.0f * y * v2);
    if (root < 0.0f) return false;

    float angle = atan2f(v2 - sqrtf(root), g * xz); // lower arc
    Vec3 horizontal = normalize((Vec3){delta.x, 0.0f, delta.z});
    *velocityOut = add(scale(horizontal, cosf(angle) * projectileSpeed),
                       (Vec3){0.0f, sinf(angle) * projectileSpeed, 0.0f});
    return true;
}
\end{lstlisting}

\Needspace{5\baselineskip}
\subsection*{21. Iterative Targeting for Moving Targets}
\addcontentsline{toc}{subsection}{21. Iterative Targeting for Moving Targets}

\begin{lstlisting}
bool iterativeTargeting(
    Vec3 shooter,
    Kinematic target,
    float projectileSpeed,
    Vec3 gravity,
    int maxIterations,
    Vec3 *velocityOut
) {
    Vec3 aimPoint = target.position;
    for (int i = 0; i < maxIterations; ++i) {
        Vec3 v;
        if (!firingSolution(shooter, aimPoint, projectileSpeed, gravity, &v)) return false;
        float t = length(sub(aimPoint, shooter)) / projectileSpeed;
        aimPoint = add(target.position, scale(target.velocity, t));
        *velocityOut = v;
    }
    return true;
}
\end{lstlisting}

\Needspace{5\baselineskip}
\subsection*{22. Jump Planning}
\addcontentsline{toc}{subsection}{22. Jump Planning}

\begin{lstlisting}
typedef struct JumpPoint {
    Vec3 takeoff;
    Vec3 landing;
    float jumpTime;
    Vec3 takeoffVelocity;
} JumpPoint;

bool calculateJump(Kinematic character, JumpPoint *jump, Vec3 gravity) {
    Vec3 delta = sub(jump->landing, jump->takeoff);
    float t = jump->jumpTime;
    if (t <= 0.0f) return false;
    jump->takeoffVelocity = sub(scale(delta, 1.0f / t), scale(gravity, 0.5f * t));
    return true;
}

SteeringOutput jumpSteering(Kinematic character, JumpPoint jump, float maxAccel) {
    Kinematic target = character;
    target.position = jump.takeoff;
    return arrive(character, target, length(jump.takeoffVelocity), maxAccel,
                  0.05f, 1.0f, 0.1f);
}
\end{lstlisting}

\Needspace{5\baselineskip}
\subsection*{23. Formation Slot Assignment}
\addcontentsline{toc}{subsection}{23. Formation Slot Assignment}

\begin{lstlisting}
typedef struct FormationSlot { int slotId; Vec3 relativePosition; float relativeOrientation; } FormationSlot;
typedef struct CharacterSlot { int characterId; int slotId; float cost; } CharacterSlot;

float slotCost(Kinematic c, FormationSlot s, Kinematic anchor) {
    Vec3 worldSlot = add(anchor.position, s.relativePosition);
    return length(sub(c.position, worldSlot));
}

void greedyAssignSlots(Kinematic *characters, int n, FormationSlot *slots, int m, int *slotForCharacter) {
    bool usedSlot[1024] = {0};
    for (int i = 0; i < n; ++i) slotForCharacter[i] = -1;

    for (int assigned = 0; assigned < n && assigned < m; ++assigned) {
        float best = INFINITY;
        int bestC = -1, bestS = -1;
        for (int c = 0; c < n; ++c) if (slotForCharacter[c] < 0) {
            for (int s = 0; s < m; ++s) if (!usedSlot[s]) {
                float cost = slotCost(characters[c], slots[s], characters[0]);
                if (cost < best) { best = cost; bestC = c; bestS = s; }
            }
        }
        slotForCharacter[bestC] = bestS;
        usedSlot[bestS] = true;
    }
}
\end{lstlisting}

\Needspace{5\baselineskip}
\subsection*{24. Formation Anchor and Drift Correction}
\addcontentsline{toc}{subsection}{24. Formation Anchor and Drift Correction}

\begin{lstlisting}
Kinematic formationAnchor(Kinematic *members, int n, FormationSlot *slots, int *slotForCharacter) {
    Vec3 sum = {0,0,0};
    int count = 0;
    for (int i = 0; i < n; ++i) {
        int s = slotForCharacter[i];
        if (s < 0) continue;
        sum = add(sum, sub(members[i].position, slots[s].relativePosition));
        ++count;
    }
    Kinematic anchor = {0};
    if (count > 0) anchor.position = scale(sum, 1.0f / count);
    return anchor;
}
\end{lstlisting}

\medskip\hrule\medskip

\clearpage
\section*{Part II -- Pathfinding Algorithms}
\addcontentsline{toc}{section}{Part II -- Pathfinding Algorithms}

\Needspace{5\baselineskip}
\subsection*{25. Graph Interfaces}
\addcontentsline{toc}{subsection}{25. Graph Interfaces}

\begin{lstlisting}
typedef struct Node Node;

typedef struct Connection {
    Node *from;
    Node *to;
    float cost;
} Connection;

typedef struct Graph {
    Connection *(*connections)(struct Graph *g, Node *node, int *countOut);
} Graph;

typedef struct NodeRecord {
    Node *node;
    Connection *connection;
    float costSoFar;
    float estimatedTotal;
} NodeRecord;
\end{lstlisting}

\Needspace{5\baselineskip}
\subsection*{26. Dijkstra Pathfinding}
\addcontentsline{toc}{subsection}{26. Dijkstra Pathfinding}

\begin{lstlisting}
Path dijkstra(Graph *graph, Node *start, Node *goal) {
    PriorityQueue open = pqCreate();
    Set closed = setCreate();
    Map records = mapCreate();

    NodeRecord startRec = { start, NULL, 0.0f, 0.0f };
    pqPush(&open, startRec, 0.0f);
    mapPut(&records, start, startRec);

    while (!pqEmpty(&open)) {
        NodeRecord current = pqPopMin(&open);
        if (current.node == goal) return reconstructPath(current, records);

        int edgeCount = 0;
        Connection *edges = graph->connections(graph, current.node, &edgeCount);
        for (int i = 0; i < edgeCount; ++i) {
            Node *end = edges[i].to;
            float endCost = current.costSoFar + edges[i].cost;
            if (setContains(&closed, end)) continue;

            NodeRecord old;
            bool seen = mapGet(&records, end, &old);
            if (!seen || endCost < old.costSoFar) {
                NodeRecord nr = { end, &edges[i], endCost, endCost };
                mapPut(&records, end, nr);
                pqUpsert(&open, nr, endCost);
            }
        }
        setAdd(&closed, current.node);
    }
    return emptyPath();
}
\end{lstlisting}

\Needspace{5\baselineskip}
\subsection*{27. A\textsuperscript{*} Pathfinding}
\addcontentsline{toc}{subsection}{27. A\textsuperscript{*} Pathfinding}

\begin{lstlisting}
typedef float (*Heuristic)(Node *a, Node *goal);

Path astar(Graph *graph, Node *start, Node *goal, Heuristic h) {
    PriorityQueue open = pqCreate();
    Set closed = setCreate();
    Map records = mapCreate();

    NodeRecord startRec = { start, NULL, 0.0f, h(start, goal) };
    pqPush(&open, startRec, startRec.estimatedTotal);
    mapPut(&records, start, startRec);

    while (!pqEmpty(&open)) {
        NodeRecord current = pqPopMin(&open);
        if (current.node == goal) return reconstructPath(current, records);

        int edgeCount = 0;
        Connection *edges = graph->connections(graph, current.node, &edgeCount);
        for (int i = 0; i < edgeCount; ++i) {
            Node *end = edges[i].to;
            float g = current.costSoFar + edges[i].cost;
            if (setContains(&closed, end)) {
                NodeRecord old = mapAt(&records, end);
                if (old.costSoFar <= g) continue;
                setRemove(&closed, end);
            }

            NodeRecord old;
            bool seen = mapGet(&records, end, &old);
            if (!seen || g < old.costSoFar) {
                NodeRecord nr = { end, &edges[i], g, g + h(end, goal) };
                mapPut(&records, end, nr);
                pqUpsert(&open, nr, nr.estimatedTotal);
            }
        }
        setAdd(&closed, current.node);
    }
    return emptyPath();
}
\end{lstlisting}

\Needspace{5\baselineskip}
\subsection*{28. Node-Array A\textsuperscript{*}}
\addcontentsline{toc}{subsection}{28. Node-Array A\textsuperscript{*}}

\begin{lstlisting}
Path astarNodeArray(Graph *graph, int nodeCount, int startId, int goalId, HeuristicById h) {
    NodeRecord *records = calloc(nodeCount, sizeof(NodeRecord));
    uint8_t *state = calloc(nodeCount, sizeof(uint8_t)); // 0 unseen, 1 open, 2 closed
    BinaryHeap open = heapCreate();

    records[startId] = (NodeRecord){ nodeById(startId), NULL, 0.0f, h(startId, goalId) };
    state[startId] = 1;
    heapPush(&open, startId, records[startId].estimatedTotal);

    while (!heapEmpty(&open)) {
        int currentId = heapPopMin(&open);
        if (currentId == goalId) return reconstructArrayPath(records, startId, goalId);
        state[currentId] = 2;

        int n = 0;
        Connection *edges = graphConnectionsById(graph, currentId, &n);
        for (int i = 0; i < n; ++i) {
            int to = idOf(edges[i].to);
            float g = records[currentId].costSoFar + edges[i].cost;
            if (state[to] == 2 && records[to].costSoFar <= g) continue;
            if (state[to] == 0 || g < records[to].costSoFar) {
                records[to] = (NodeRecord){ edges[i].to, &edges[i], g, g + h(to, goalId) };
                state[to] = 1;
                heapUpsert(&open, to, records[to].estimatedTotal);
            }
        }
    }
    return emptyPath();
}
\end{lstlisting}

\Needspace{5\baselineskip}
\subsection*{29. Path Smoothing}
\addcontentsline{toc}{subsection}{29. Path Smoothing}

\begin{lstlisting}
Path smoothPath(Path raw, bool (*hasLineOfSight)(Vec3 a, Vec3 b)) {
    Path out = pathCreate();
    int anchor = 0;
    pathAppend(&out, raw.points[0]);

    while (anchor < raw.count - 1) {
        int farthest = anchor + 1;
        for (int i = raw.count - 1; i > anchor; --i) {
            if (hasLineOfSight(raw.points[anchor], raw.points[i])) {
                farthest = i;
                break;
            }
        }
        pathAppend(&out, raw.points[farthest]);
        anchor = farthest;
    }
    return out;
}
\end{lstlisting}

\Needspace{5\baselineskip}
\subsection*{30. Hierarchical Pathfinding}
\addcontentsline{toc}{subsection}{30. Hierarchical Pathfinding}

\begin{lstlisting}
Path hierarchicalPathfind(Hierarchy *h, Node *start, Node *goal, Heuristic heuristic) {
    int top = h->levelCount - 1;
    Node *abstractStart = h->abstractNode(start, top);
    Node *abstractGoal = h->abstractNode(goal, top);
    Path abstractPath = astar(h->graph[top], abstractStart, abstractGoal, heuristic);

    Path refined = emptyPath();
    Node *localStart = start;
    for (int i = 1; i < abstractPath.count; ++i) {
        Node *entry = h->entryNode(abstractPath.nodes[i-1], abstractPath.nodes[i]);
        Path segment = astar(h->graph[0], localStart, entry, heuristic);
        pathConcat(&refined, segment);
        localStart = entry;
    }
    pathConcat(&refined, astar(h->graph[0], localStart, goal, heuristic));
    return refined;
}
\end{lstlisting}

\Needspace{5\baselineskip}
\subsection*{31. Continuous-Time A\textsuperscript{*}}
\addcontentsline{toc}{subsection}{31. Continuous-Time A\textsuperscript{*}}

\begin{lstlisting}
typedef struct TimeNode { Vec3 position; float time; } TimeNode;

Path continuousTimeAStar(TimeGraph *graph, TimeNode start, TimeNode goal, HeuristicTime h) {
    PriorityQueue open = pqCreate();
    Map records = mapCreate();
    pqPushTime(&open, start, h(start, goal));

    while (!pqEmpty(&open)) {
        TimeNode current = pqPopTime(&open);
        if (timeNodeGoalReached(current, goal)) return reconstructTimePath(current, records);

        int n = 0;
        TimeConnection *moves = graph->connections(graph, current, &n);
        for (int i = 0; i < n; ++i) {
            TimeNode next = moves[i].to;
            if (!graph->isCollisionFree(graph, current, next)) continue;
            float g = recordCost(records, current) + moves[i].cost;
            if (recordBetter(records, next, g)) continue;
            storeRecord(records, next, current, moves[i], g, g + h(next, goal));
            pqUpsertTime(&open, next, g + h(next, goal));
        }
    }
    return emptyPath();
}
\end{lstlisting}

\Needspace{5\baselineskip}
\subsection*{32. Movement Planning over Action Graphs}
\addcontentsline{toc}{subsection}{32. Movement Planning over Action Graphs}

\begin{lstlisting}
typedef struct MovementAction {
    bool (*precondition)(WorldState s);
    WorldState (*apply)(WorldState s);
    float cost;
} MovementAction;

Plan movementPlan(WorldState start, bool (*goal)(WorldState), MovementAction *actions, int actionCount) {
    SearchNode root = { start, NULL, NULL, 0.0f };
    PriorityQueue open = pqCreate();
    pqPush(&open, root, 0.0f);

    while (!pqEmpty(&open)) {
        SearchNode cur = pqPopMin(&open);
        if (goal(cur.state)) return reconstructPlan(cur);
        for (int i = 0; i < actionCount; ++i) {
            if (!actions[i].precondition(cur.state)) continue;
            WorldState next = actions[i].apply(cur.state);
            SearchNode child = { next, &cur, &actions[i], cur.cost + actions[i].cost };
            pqPush(&open, child, child.cost);
        }
    }
    return emptyPlan();
}
\end{lstlisting}

\medskip\hrule\medskip

\clearpage
\section*{Part III -- Decision-Making Algorithms}
\addcontentsline{toc}{section}{Part III -- Decision-Making Algorithms}

\Needspace{5\baselineskip}
\subsection*{33. Decision Tree Evaluation}
\addcontentsline{toc}{subsection}{33. Decision Tree Evaluation}

\begin{lstlisting}
typedef struct DecisionNode DecisionNode;

typedef struct DecisionNode {
    bool isLeaf;
    Action action;
    bool (*test)(GameState *s, struct DecisionNode *self);
    DecisionNode *trueBranch;
    DecisionNode *falseBranch;
} DecisionNode;

Action decide(DecisionNode *node, GameState *state) {
    while (!node->isLeaf) {
        node = node->test(state, node) ? node->trueBranch : node->falseBranch;
    }
    return node->action;
}
\end{lstlisting}

\Needspace{5\baselineskip}
\subsection*{34. Random Decision Node}
\addcontentsline{toc}{subsection}{34. Random Decision Node}

\begin{lstlisting}
Action randomDecision(Action *actions, float *weights, int count) {
    float total = 0.0f;
    for (int i = 0; i < count; ++i) total += weights[i];
    float r = random01() * total;
    for (int i = 0; i < count; ++i) {
        r -= weights[i];
        if (r <= 0.0f) return actions[i];
    }
    return actions[count - 1];
}
\end{lstlisting}

\Needspace{5\baselineskip}
\subsection*{35. Finite State Machine Update}
\addcontentsline{toc}{subsection}{35. Finite State Machine Update}

\begin{lstlisting}
typedef struct Transition {
    bool (*triggered)(GameState *s, void *ctx);
    int targetState;
    Action *actions;
    int actionCount;
    void *ctx;
} Transition;

typedef struct State {
    Action *actions;
    int actionCount;
    Transition *transitions;
    int transitionCount;
    Action *entryActions;
    int entryCount;
    Action *exitActions;
    int exitCount;
} State;

void updateFSM(StateMachine *fsm, GameState *state) {
    State *cur = &fsm->states[fsm->current];
    for (int i = 0; i < cur->transitionCount; ++i) {
        Transition *t = &cur->transitions[i];
        if (t->triggered(state, t->ctx)) {
            executeActions(cur->exitActions, cur->exitCount, state);
            executeActions(t->actions, t->actionCount, state);
            fsm->current = t->targetState;
            executeActions(fsm->states[fsm->current].entryActions,
                           fsm->states[fsm->current].entryCount, state);
            return;
        }
    }
    executeActions(cur->actions, cur->actionCount, state);
}
\end{lstlisting}

\Needspace{5\baselineskip}
\subsection*{36. Hierarchical State Machine Update}
\addcontentsline{toc}{subsection}{36. Hierarchical State Machine Update}

\begin{lstlisting}
UpdateResult updateHSM(HSMState *state, GameState *world) {
    for (int i = 0; i < state->transitionCount; ++i) {
        Transition *t = &state->transitions[i];
        if (t->triggered(world, t->ctx)) return (UpdateResult){ .transition = t, .handled = true };
    }
    if (state->child) {
        UpdateResult r = updateHSM(state->child, world);
        if (r.handled) return r;
    }
    executeActions(state->actions, state->actionCount, world);
    return (UpdateResult){0};
}
\end{lstlisting}

\Needspace{5\baselineskip}
\subsection*{37. Behavior Tree Selector and Sequence}
\addcontentsline{toc}{subsection}{37. Behavior Tree Selector and Sequence}

\begin{lstlisting}
typedef enum BTStatus { BT_SUCCESS, BT_FAILURE, BT_RUNNING } BTStatus;
typedef BTStatus (*BTNodeRun)(void *node, GameState *s);

BTStatus selectorRun(BTNode **children, int count, GameState *s) {
    for (int i = 0; i < count; ++i) {
        BTStatus r = children[i]->run(children[i], s);
        if (r != BT_FAILURE) return r;
    }
    return BT_FAILURE;
}

BTStatus sequenceRun(BTNode **children, int count, GameState *s) {
    for (int i = 0; i < count; ++i) {
        BTStatus r = children[i]->run(children[i], s);
        if (r != BT_SUCCESS) return r;
    }
    return BT_SUCCESS;
}
\end{lstlisting}

\Needspace{5\baselineskip}
\subsection*{38. Behavior Tree Random and Non-Deterministic Composites}
\addcontentsline{toc}{subsection}{38. Behavior Tree Random and Non-Deterministic Composites}

\begin{lstlisting}
BTStatus randomSelectorRun(BTNode **children, float *weights, int count, GameState *s) {
    int start = weightedRandomIndex(weights, count);
    for (int k = 0; k < count; ++k) {
        int i = (start + k) % count;
        BTStatus r = children[i]->run(children[i], s);
        if (r != BT_FAILURE) return r;
    }
    return BT_FAILURE;
}

BTStatus shuffledSequenceRun(BTNode **children, int count, GameState *s) {
    int order[MAX_CHILDREN];
    randomPermutation(order, count);
    for (int k = 0; k < count; ++k) {
        BTStatus r = children[order[k]]->run(children[order[k]], s);
        if (r != BT_SUCCESS) return r;
    }
    return BT_SUCCESS;
}
\end{lstlisting}

\Needspace{5\baselineskip}
\subsection*{39. Behavior Tree Decorators}
\addcontentsline{toc}{subsection}{39. Behavior Tree Decorators}

\begin{lstlisting}
BTStatus inverterRun(BTNode *child, GameState *s) {
    BTStatus r = child->run(child, s);
    if (r == BT_SUCCESS) return BT_FAILURE;
    if (r == BT_FAILURE) return BT_SUCCESS;
    return BT_RUNNING;
}

BTStatus limitRun(LimitNode *node, GameState *s) {
    if (node->used >= node->limit) return BT_FAILURE;
    BTStatus r = node->child->run(node->child, s);
    if (r == BT_SUCCESS || r == BT_FAILURE) node->used++;
    return r;
}

BTStatus untilFailRun(BTNode *child, GameState *s) {
    while (true) {
        BTStatus r = child->run(child, s);
        if (r == BT_FAILURE) return BT_SUCCESS;
        if (r == BT_RUNNING) return BT_RUNNING;
    }
}
\end{lstlisting}

\Needspace{5\baselineskip}
\subsection*{40. Parallel Behavior Tree}
\addcontentsline{toc}{subsection}{40. Parallel Behavior Tree}

\begin{lstlisting}
BTStatus parallelRun(BTNode **children, int count, int successThreshold, int failureThreshold, GameState *s) {
    int successes = 0, failures = 0;
    for (int i = 0; i < count; ++i) {
        BTStatus r = children[i]->run(children[i], s);
        if (r == BT_SUCCESS) successes++;
        else if (r == BT_FAILURE) failures++;
    }
    if (successes >= successThreshold) return BT_SUCCESS;
    if (failures >= failureThreshold) return BT_FAILURE;
    return BT_RUNNING;
}
\end{lstlisting}

\Needspace{5\baselineskip}
\subsection*{41. Blackboard Behavior Tree Lookup}
\addcontentsline{toc}{subsection}{41. Blackboard Behavior Tree Lookup}

\begin{lstlisting}
BTStatus blackboardTaskRun(Blackboard *bb, const char *key, BTNode *child, GameState *s) {
    void *value = blackboardGet(bb, key);
    if (value == NULL) return BT_FAILURE;
    pushTaskContext(s, value);
    BTStatus r = child->run(child, s);
    popTaskContext(s);
    return r;
}
\end{lstlisting}

\Needspace{5\baselineskip}
\subsection*{42. Fuzzy Decision Making}
\addcontentsline{toc}{subsection}{42. Fuzzy Decision Making}

\begin{lstlisting}
typedef float (*Membership)(float x);

typedef struct FuzzyRule {
    Membership inputSet;
    float (*outputValue)(GameState *s);
    float weight;
} FuzzyRule;

float fuzzyDecision(FuzzyRule *rules, int count, float observed, GameState *s) {
    float weighted = 0.0f;
    float total = 0.0f;
    for (int i = 0; i < count; ++i) {
        float dom = rules[i].inputSet(observed) * rules[i].weight;
        weighted += dom * rules[i].outputValue(s);
        total += dom;
    }
    return (total > 0.0f) ? weighted / total : 0.0f;
}
\end{lstlisting}

\Needspace{5\baselineskip}
\subsection*{43. Fuzzy State Machine}
\addcontentsline{toc}{subsection}{43. Fuzzy State Machine}

\begin{lstlisting}
void updateFuzzyStateMachine(FuzzyFSM *fsm, GameState *s) {
    for (int i = 0; i < fsm->stateCount; ++i) fsm->nextDom[i] = 0.0f;

    for (int i = 0; i < fsm->stateCount; ++i) {
        float dom = fsm->dom[i];
        if (dom <= 0.0f) continue;
        executeWeightedActions(fsm->states[i].actions, dom, s);
        for (int t = 0; t < fsm->states[i].transitionCount; ++t) {
            FuzzyTransition *tr = &fsm->states[i].transitions[t];
            float triggered = tr->membership(s) * dom;
            fsm->nextDom[tr->target] = fmaxf(fsm->nextDom[tr->target], triggered);
        }
    }
    normalizeArray(fsm->nextDom, fsm->stateCount);
    memcpy(fsm->dom, fsm->nextDom, sizeof(float) * fsm->stateCount);
}
\end{lstlisting}

\Needspace{5\baselineskip}
\subsection*{44. Markov State Machine}
\addcontentsline{toc}{subsection}{44. Markov State Machine}

\begin{lstlisting}
int markovNextState(float *transitionRow, int count) {
    float r = random01();
    float cumulative = 0.0f;
    for (int i = 0; i < count; ++i) {
        cumulative += transitionRow[i];
        if (r <= cumulative) return i;
    }
    return count - 1;
}

void updateMarkovStateMachine(MarkovFSM *fsm, GameState *s) {
    executeActions(fsm->states[fsm->current].actions, fsm->states[fsm->current].actionCount, s);
    fsm->current = markovNextState(fsm->transition[fsm->current], fsm->stateCount);
}
\end{lstlisting}

\Needspace{5\baselineskip}
\subsection*{45. Goal-Oriented Behavior Selection}
\addcontentsline{toc}{subsection}{45. Goal-Oriented Behavior Selection}

\begin{lstlisting}
typedef struct Goal { float value; float insistence; } GoalValue;
typedef struct ActionModel { float duration; float *goalDelta; bool (*canRun)(GameState *s); } ActionModel;

float discontentment(GoalValue *goals, int count) {
    float d = 0.0f;
    for (int i = 0; i < count; ++i) d += goals[i].insistence * goals[i].value * goals[i].value;
    return d;
}

ActionModel *chooseGoalAction(ActionModel *actions, int actionCount, GoalValue *goals, int goalCount, GameState *s) {
    float best = INFINITY;
    ActionModel *bestAction = NULL;
    for (int a = 0; a < actionCount; ++a) {
        if (!actions[a].canRun(s)) continue;
        GoalValue simulated[MAX_GOALS];
        memcpy(simulated, goals, sizeof(GoalValue) * goalCount);
        for (int g = 0; g < goalCount; ++g) simulated[g].value += actions[a].goalDelta[g];
        float score = discontentment(simulated, goalCount);
        if (score < best) { best = score; bestAction = &actions[a]; }
    }
    return bestAction;
}
\end{lstlisting}

\Needspace{5\baselineskip}
\subsection*{46. Depth-Limited GOAP Planning}
\addcontentsline{toc}{subsection}{46. Depth-Limited GOAP Planning}

\begin{lstlisting}
Plan planActionDFS(WorldModel world, ActionModel *actions, int actionCount, int maxDepth) {
    float bestScore = worldDiscontentment(world);
    Plan best = emptyPlan();
    dfsPlan(world, actions, actionCount, maxDepth, emptyPlan(), &best, &bestScore);
    return best;
}

void dfsPlan(WorldModel world, ActionModel *actions, int n, int depth, Plan prefix, Plan *best, float *bestScore) {
    float score = worldDiscontentment(world);
    if (score < *bestScore) { *bestScore = score; *best = prefix; }
    if (depth == 0) return;

    for (int i = 0; i < n; ++i) {
        if (!actionCanRun(actions[i], world)) continue;
        WorldModel next = applyAction(world, actions[i]);
        Plan extended = appendPlan(prefix, actions[i]);
        dfsPlan(next, actions, n, depth - 1, extended, best, bestScore);
    }
}
\end{lstlisting}

\Needspace{5\baselineskip}
\subsection*{47. GOAP with Iterative Deepening A\textsuperscript{*}}
\addcontentsline{toc}{subsection}{47. GOAP with Iterative Deepening A\textsuperscript{*}}

\begin{lstlisting}
Plan goapIDAStar(WorldModel start, ActionModel *actions, int actionCount, float (*heuristic)(WorldModel)) {
    float bound = heuristic(start);
    Plan path = emptyPlan();
    while (true) {
        float t = idaSearch(start, actions, actionCount, 0.0f, bound, &path);
        if (t == FOUND) return path;
        if (t == INFINITY) return emptyPlan();
        bound = t;
    }
}
\end{lstlisting}

\Needspace{5\baselineskip}
\subsection*{48. Rule-Based System Iteration}
\addcontentsline{toc}{subsection}{48. Rule-Based System Iteration}

\begin{lstlisting}
typedef struct Rule {
    bool (*matches)(Database *db, Bindings *out);
    void (*act)(Bindings b, Database *db, GameState *s);
    int priority;
} Rule;

void ruleBasedIteration(Rule *rules, int count, Database *db, GameState *s) {
    Candidate candidates[MAX_RULES];
    int n = 0;
    for (int i = 0; i < count; ++i) {
        Bindings b;
        if (rules[i].matches(db, &b)) candidates[n++] = (Candidate){ &rules[i], b };
    }
    if (n == 0) return;
    Candidate chosen = arbitrateRule(candidates, n);
    chosen.rule->act(chosen.bindings, db, s);
}
\end{lstlisting}

\Needspace{5\baselineskip}
\subsection*{49. Pattern Matching with Bindings}
\addcontentsline{toc}{subsection}{49. Pattern Matching with Bindings}

\begin{lstlisting}
bool matchDatum(Pattern p, Datum d, Bindings *b) {
    if (p.kind != d.kind) return false;
    for (int i = 0; i < p.fieldCount; ++i) {
        if (p.fields[i].isVariable) {
            if (!bindOrCheck(b, p.fields[i].name, d.fields[i])) return false;
        } else if (!valueEquals(p.fields[i].literal, d.fields[i])) {
            return false;
        }
    }
    return true;
}
\end{lstlisting}

\Needspace{5\baselineskip}
\subsection*{50. Rete-Style Fact Propagation}
\addcontentsline{toc}{subsection}{50. Rete-Style Fact Propagation}

\begin{lstlisting}
void reteAddFact(ReteNetwork *net, Datum fact) {
    for (int i = 0; i < net->patternCount; ++i) {
        PatternNode *p = &net->patterns[i];
        Bindings b;
        if (matchDatum(p->pattern, fact, &b)) {
            Match m = { fact, b };
            patternStore(p, m);
            forEachOutput(p, propagateMatch, &m);
        }
    }
}

void reteRemoveFact(ReteNetwork *net, Datum fact) {
    for (int i = 0; i < net->patternCount; ++i) {
        if (patternContains(&net->patterns[i], fact)) {
            patternRemove(&net->patterns[i], fact);
            propagateRemoval(&net->patterns[i], fact);
        }
    }
}
\end{lstlisting}

\Needspace{5\baselineskip}
\subsection*{51. Blackboard Iteration}
\addcontentsline{toc}{subsection}{51. Blackboard Iteration}

\begin{lstlisting}
void blackboardIteration(Blackboard *bb, Expert *experts, int count, GameState *s) {
    float best = -INFINITY;
    Expert *chosen = NULL;
    for (int i = 0; i < count; ++i) {
        float insistence = experts[i].getInsistence(bb, s);
        if (insistence > best) { best = insistence; chosen = &experts[i]; }
    }
    if (chosen != NULL) chosen->run(bb, s);
}
\end{lstlisting}

\Needspace{5\baselineskip}
\subsection*{52. Action Execution Manager}
\addcontentsline{toc}{subsection}{52. Action Execution Manager}

\begin{lstlisting}
bool scheduleAction(ActionManager *mgr, Action *a) {
    for (int i = 0; i < mgr->activeCount; ++i) {
        Action *other = mgr->active[i];
        if (!a->canDoBoth(a, other)) {
            if (!other->interrupt(other)) return false;
            removeActive(mgr, i--);
        }
    }
    mgr->active[mgr->activeCount++] = a;
    return true;
}

void executeActionsManaged(ActionManager *mgr, GameState *s, float dt) {
    for (int i = 0; i < mgr->activeCount; ++i) {
        mgr->active[i]->execute(mgr->active[i], s, dt);
        if (mgr->active[i]->isComplete(mgr->active[i])) removeActive(mgr, i--);
    }
}
\end{lstlisting}

\medskip\hrule\medskip

\clearpage
\section*{Part IV -- Tactical and Strategic AI}
\addcontentsline{toc}{section}{Part IV -- Tactical and Strategic AI}

\Needspace{5\baselineskip}
\subsection*{53. Cover Quality Evaluation}
\addcontentsline{toc}{subsection}{53. Cover Quality Evaluation}

\begin{lstlisting}
float coverQuality(Vec3 cover, Vec3 enemy, Vec3 self, World *world) {
    float visibilityPenalty = worldLineOfSight(world, enemy, cover) ? 1.0f : 0.0f;
    float distanceScore = 1.0f / (1.0f + length(sub(self, cover)));
    float attackScore = worldLineOfSight(world, cover, enemy) ? 1.0f : 0.25f;
    return 2.0f * attackScore + distanceScore - 3.0f * visibilityPenalty;
}
\end{lstlisting}

\Needspace{5\baselineskip}
\subsection*{54. Waypoint Condensation}
\addcontentsline{toc}{subsection}{54. Waypoint Condensation}

\begin{lstlisting}
void condenseWaypoints(WaypointList *list, float mergeDistance) {
    listSortByImportanceDescending(list);
    for (Waypoint *w = listFirst(list); w != NULL; w = listNext(list, w)) {
        Waypoint **nearby = listNearby(list, w->position, mergeDistance);
        for (int i = 0; nearby[i] != NULL; ++i) {
            if (nearby[i] == w) continue;
            if (nearby[i]->importance <= w->importance) {
                mergeWaypointData(w, nearby[i]);
                listRemove(list, nearby[i]);
            }
        }
    }
}
\end{lstlisting}

\Needspace{5\baselineskip}
\subsection*{55. Influence Map Flooding}
\addcontentsline{toc}{subsection}{55. Influence Map Flooding}

\begin{lstlisting}
void mapFloodDijkstra(Map *map, Location *sources, float *sourceStrengths, int sourceCount) {
    PriorityQueue open = pqCreate();
    for (int i = 0; i < sourceCount; ++i) {
        sources[i].influence = sourceStrengths[i];
        pqPush(&open, sources[i], -sourceStrengths[i]);
    }

    while (!pqEmpty(&open)) {
        Location cur = pqPopMin(&open);
        int n = 0;
        Location *neighbors = map->neighbors(map, cur, &n);
        for (int i = 0; i < n; ++i) {
            float candidate = cur.influence * attenuation(cur, neighbors[i]);
            if (candidate > neighbors[i].influence) {
                neighbors[i].influence = candidate;
                pqUpsert(&open, neighbors[i], -candidate);
            }
        }
    }
}
\end{lstlisting}

\Needspace{5\baselineskip}
\subsection*{56. Convolution Filter}
\addcontentsline{toc}{subsection}{56. Convolution Filter}

\begin{lstlisting}
void convolve(float **src, float **dst, int width, int height, float **kernel, int radius) {
    for (int y = 0; y < height; ++y) {
        for (int x = 0; x < width; ++x) {
            float sum = 0.0f;
            for (int ky = -radius; ky <= radius; ++ky) {
                for (int kx = -radius; kx <= radius; ++kx) {
                    int sx = clampInt(x + kx, 0, width - 1);
                    int sy = clampInt(y + ky, 0, height - 1);
                    sum += src[sy][sx] * kernel[ky + radius][kx + radius];
                }
            }
            dst[y][x] = sum;
        }
    }
}
\end{lstlisting}

\Needspace{5\baselineskip}
\subsection*{57. Separable Convolution}
\addcontentsline{toc}{subsection}{57. Separable Convolution}

\begin{lstlisting}
void separableConvolve(float **src, float **dst, int width, int height, float *kernel, int radius) {
    float **tmp = allocateGrid(width, height);
    for (int y = 0; y < height; ++y)
        for (int x = 0; x < width; ++x) {
            float sum = 0.0f;
            for (int k = -radius; k <= radius; ++k)
                sum += src[y][clampInt(x+k,0,width-1)] * kernel[k+radius];
            tmp[y][x] = sum;
        }

    for (int y = 0; y < height; ++y)
        for (int x = 0; x < width; ++x) {
            float sum = 0.0f;
            for (int k = -radius; k <= radius; ++k)
                sum += tmp[clampInt(y+k,0,height-1)][x] * kernel[k+radius];
            dst[y][x] = sum;
        }
}
\end{lstlisting}

\Needspace{5\baselineskip}
\subsection*{58. Cellular Automata Update}
\addcontentsline{toc}{subsection}{58. Cellular Automata Update}

\begin{lstlisting}
void cellularAutomataStep(Grid *grid, Grid *next, CellRule rule) {
    for (int y = 0; y < grid->height; ++y) {
        for (int x = 0; x < grid->width; ++x) {
            Neighborhood n = collectNeighborhood(grid, x, y);
            next->cell[y][x] = rule(grid->cell[y][x], n);
        }
    }
    swapGridContents(grid, next);
}
\end{lstlisting}

\Needspace{5\baselineskip}
\subsection*{59. Tactical Path Cost Blending}
\addcontentsline{toc}{subsection}{59. Tactical Path Cost Blending}

\begin{lstlisting}
float tacticalEdgeCost(Connection c, TacticalMap *maps, float *weights, int count) {
    float cost = c.cost;
    Vec3 p = c.to->position;
    for (int i = 0; i < count; ++i) cost += weights[i] * sampleTacticalMap(&maps[i], p);
    return fmaxf(cost, 0.001f);
}
\end{lstlisting}

\medskip\hrule\medskip

\clearpage
\section*{Part V -- Learning Algorithms}
\addcontentsline{toc}{section}{Part V -- Learning Algorithms}

\Needspace{5\baselineskip}
\subsection*{60. Parameter Optimization Driver}
\addcontentsline{toc}{subsection}{60. Parameter Optimization Driver}

\begin{lstlisting}
Parameters optimizeParameters(Parameters start, float (*fitness)(Parameters), OptimizerStep step, int iterations) {
    Parameters p = start;
    float best = fitness(p);
    for (int i = 0; i < iterations; ++i) {
        Parameters q = step(p, i);
        float score = fitness(q);
        if (score > best) { p = q; best = score; }
    }
    return p;
}
\end{lstlisting}

\Needspace{5\baselineskip}
\subsection*{61. Hill Climbing}
\addcontentsline{toc}{subsection}{61. Hill Climbing}

\begin{lstlisting}
Parameters hillClimb(Parameters p, float stepSize, float (*fitness)(Parameters), int iterations) {
    float best = fitness(p);
    for (int i = 0; i < iterations; ++i) {
        Parameters candidate = p;
        int dim = randomInt(0, p.count - 1);
        candidate.value[dim] += randomSigned() * stepSize;
        float score = fitness(candidate);
        if (score > best) { p = candidate; best = score; }
    }
    return p;
}
\end{lstlisting}

\Needspace{5\baselineskip}
\subsection*{62. Simulated Annealing}
\addcontentsline{toc}{subsection}{62. Simulated Annealing}

\begin{lstlisting}
Parameters anneal(Parameters p, float (*fitness)(Parameters), int iterations, float startTemp, float cooling) {
    float current = fitness(p);
    float temp = startTemp;
    for (int i = 0; i < iterations; ++i) {
        Parameters q = randomNeighbor(p);
        float score = fitness(q);
        float delta = score - current;
        if (delta > 0.0f || random01() < expf(delta / temp)) {
            p = q;
            current = score;
        }
        temp *= cooling;
    }
    return p;
}
\end{lstlisting}

\Needspace{5\baselineskip}
\subsection*{63. N-Gram Action Prediction}
\addcontentsline{toc}{subsection}{63. N-Gram Action Prediction}

\begin{lstlisting}
typedef struct NGramPredictor { Map counts; int n; } NGramPredictor;

void ngramRegister(NGramPredictor *p, ActionSymbol *sequence, int length) {
    for (int i = 0; i + p->n <= length; ++i) {
        Key k = makeKey(&sequence[i], p->n - 1);
        ActionSymbol next = sequence[i + p->n - 1];
        incrementCount(&p->counts, k, next);
    }
}

ActionSymbol ngramPredict(NGramPredictor *p, ActionSymbol *history) {
    Key k = makeKey(history, p->n - 1);
    return mostFrequentNext(&p->counts, k);
}
\end{lstlisting}

\Needspace{5\baselineskip}
\subsection*{64. Hierarchical N-Gram Prediction}
\addcontentsline{toc}{subsection}{64. Hierarchical N-Gram Prediction}

\begin{lstlisting}
ActionSymbol hierarchicalNgramPredict(NGramPredictor *predictors, int levels, ActionSymbol *history) {
    for (int level = levels - 1; level >= 0; --level) {
        ActionSymbol a;
        if (tryNgramPredict(&predictors[level], history, &a)) return a;
    }
    return defaultAction();
}
\end{lstlisting}

\Needspace{5\baselineskip}
\subsection*{65. Naive Bayes Classifier}
\addcontentsline{toc}{subsection}{65. Naive Bayes Classifier}

\begin{lstlisting}
void nbUpdate(NaiveBayes *nb, FeatureVector x, int classId) {
    nb->classCount[classId]++;
    for (int f = 0; f < x.count; ++f) incrementFeatureCount(nb, classId, f, x.value[f]);
}

int nbPredict(NaiveBayes *nb, FeatureVector x) {
    float best = -INFINITY;
    int bestClass = -1;
    for (int c = 0; c < nb->classTotal; ++c) {
        float logp = logf(prior(nb, c));
        for (int f = 0; f < x.count; ++f) logp += logf(likelihood(nb, c, f, x.value[f]));
        if (logp > best) { best = logp; bestClass = c; }
    }
    return bestClass;
}
\end{lstlisting}

\Needspace{5\baselineskip}
\subsection*{66. ID3 Decision Tree Learning}
\addcontentsline{toc}{subsection}{66. ID3 Decision Tree Learning}

\begin{lstlisting}
DecisionNode *makeID3Tree(Example *examples, int exampleCount, Attribute *attributes, int attributeCount) {
    if (allSameClass(examples, exampleCount)) return leafNode(examples[0].classId);
    if (attributeCount == 0) return leafNode(majorityClass(examples, exampleCount));

    int bestAttr = chooseMaxInformationGain(examples, exampleCount, attributes, attributeCount);
    DecisionNode *node = testNode(attributes[bestAttr]);

    forEachValue(attributes[bestAttr], value) {
        Example subset[MAX_EXAMPLES];
        int n = filterExamples(examples, exampleCount, attributes[bestAttr], value, subset);
        if (n == 0) node->branch[value] = leafNode(majorityClass(examples, exampleCount));
        else node->branch[value] = makeID3Tree(subset, n, removeAttr(attributes, attributeCount, bestAttr), attributeCount - 1);
    }
    return node;
}
\end{lstlisting}

\begin{lstlisting}
float entropy(ClassCounts counts) {
    float total = countTotal(counts);
    float h = 0.0f;
    for (int c = 0; c < counts.count; ++c) {
        float p = counts.value[c] / total;
        if (p > 0.0f) h -= p * log2f(p);
    }
    return h;
}
\end{lstlisting}

\Needspace{5\baselineskip}
\subsection*{67. Continuous Attribute Split Selection}
\addcontentsline{toc}{subsection}{67. Continuous Attribute Split Selection}

\begin{lstlisting}
float bestContinuousSplit(Example *examples, int n, int attr, float *thresholdOut) {
    sortExamplesByAttribute(examples, n, attr);
    float bestGain = -INFINITY;
    for (int i = 1; i < n; ++i) {
        float threshold = 0.5f * (examples[i-1].value[attr] + examples[i].value[attr]);
        Split s = splitByThreshold(examples, n, attr, threshold);
        float gain = informationGain(s);
        if (gain > bestGain) { bestGain = gain; *thresholdOut = threshold; }
    }
    return bestGain;
}
\end{lstlisting}

\Needspace{5\baselineskip}
\subsection*{68. Q-Learning}
\addcontentsline{toc}{subsection}{68. Q-Learning}

\begin{lstlisting}
void qLearningTrain(ReinforcementProblem *problem, QTable *q, int episodes, float alpha, float gamma, float rho) {
    for (int e = 0; e < episodes; ++e) {
        State s = problem->randomState(problem);
        while (!problem->terminal(problem, s)) {
            Action a = (random01() < rho) ? problem->randomAction(problem, s) : qBestAction(q, s);
            StepResult r = problem->takeAction(problem, s, a);
            float old = qGet(q, s, a);
            float future = qBestValue(q, r.nextState);
            float updated = (1.0f - alpha) * old + alpha * (r.reward + gamma * future);
            qSet(q, s, a, updated);
            s = r.nextState;
        }
    }
}
\end{lstlisting}

\Needspace{5\baselineskip}
\subsection*{69. Multilayer Perceptron Feedforward and Backpropagation}
\addcontentsline{toc}{subsection}{69. Multilayer Perceptron Feedforward and Backpropagation}

\begin{lstlisting}
void mlpFeedForward(MLP *net, float *input) {
    copy(net->layer[0].value, input, net->layer[0].count);
    for (int l = 1; l < net->layerCount; ++l) {
        for (int j = 0; j < net->layer[l].count; ++j) {
            float sum = net->layer[l].bias[j];
            for (int i = 0; i < net->layer[l-1].count; ++i)
                sum += net->weight[l][i][j] * net->layer[l-1].value[i];
            net->layer[l].value[j] = sigmoid(sum);
        }
    }
}

void mlpBackprop(MLP *net, float *target, float rate) {
    int last = net->layerCount - 1;
    for (int j = 0; j < net->layer[last].count; ++j) {
        float out = net->layer[last].value[j];
        net->layer[last].delta[j] = (target[j] - out) * out * (1.0f - out);
    }
    for (int l = last - 1; l > 0; --l) {
        for (int i = 0; i < net->layer[l].count; ++i) {
            float sum = 0.0f;
            for (int j = 0; j < net->layer[l+1].count; ++j) sum += net->weight[l+1][i][j] * net->layer[l+1].delta[j];
            float out = net->layer[l].value[i];
            net->layer[l].delta[i] = sum * out * (1.0f - out);
        }
    }
    for (int l = 1; l < net->layerCount; ++l)
        for (int i = 0; i < net->layer[l-1].count; ++i)
            for (int j = 0; j < net->layer[l].count; ++j)
                net->weight[l][i][j] += rate * net->layer[l].delta[j] * net->layer[l-1].value[i];
}
\end{lstlisting}

\medskip\hrule\medskip

\clearpage
\section*{Part VI -- Procedural Content Generation Algorithms}
\addcontentsline{toc}{section}{Part VI -- Procedural Content Generation Algorithms}

\Needspace{5\baselineskip}
\subsection*{70. Numeric Mixing for Game Seeds}
\addcontentsline{toc}{subsection}{70. Numeric Mixing for Game Seeds}

\begin{lstlisting}
uint64_t mix64(uint64_t x) {
    x ^= x >> 33;
    x *= 0xff51afd7ed558ccdULL;
    x ^= x >> 33;
    x *= 0xc4ceb9fe1a85ec53ULL;
    x ^= x >> 33;
    return x;
}
\end{lstlisting}

\Needspace{5\baselineskip}
\subsection*{71. Halton Sequence}
\addcontentsline{toc}{subsection}{71. Halton Sequence}

\begin{lstlisting}
float halton1D(int index, int base) {
    float result = 0.0f;
    float fraction = 1.0f / base;
    while (index > 0) {
        result += fraction * (index % base);
        index /= base;
        fraction /= base;
    }
    return result;
}

Vec3 halton2D(int index) {
    return (Vec3){ halton1D(index, 2), 0.0f, halton1D(index, 3) };
}
\end{lstlisting}

\Needspace{5\baselineskip}
\subsection*{72. Phyllotaxic Placement}
\addcontentsline{toc}{subsection}{72. Phyllotaxic Placement}

\begin{lstlisting}
Vec3 phyllotaxicPoint(int i, float scale) {
    const float goldenAngle = 2.39996323f;
    float r = scale * sqrtf((float)i);
    float theta = i * goldenAngle;
    return (Vec3){ r * cosf(theta), 0.0f, r * sinf(theta) };
}
\end{lstlisting}

\Needspace{5\baselineskip}
\subsection*{73. Poisson Disk Sampling}
\addcontentsline{toc}{subsection}{73. Poisson Disk Sampling}

\begin{lstlisting}
PointList poissonDisk(Rect bounds, float radius, int attempts) {
    Grid grid = gridForRadius(bounds, radius);
    ActiveList active = activeCreate();
    PointList points = pointListCreate();

    Vec3 first = randomPoint(bounds);
    addSample(&grid, &points, &active, first);

    while (!activeEmpty(&active)) {
        Vec3 p = activeRandom(&active);
        bool found = false;
        for (int k = 0; k < attempts; ++k) {
            Vec3 q = randomPointInAnnulus(p, radius, 2.0f * radius);
            if (inside(bounds, q) && noNeighborWithin(&grid, q, radius)) {
                addSample(&grid, &points, &active, q);
                found = true;
                break;
            }
        }
        if (!found) activeRemove(&active, p);
    }
    return points;
}
\end{lstlisting}

\Needspace{5\baselineskip}
\subsection*{74. Lindenmayer System Expansion}
\addcontentsline{toc}{subsection}{74. Lindenmayer System Expansion}

\begin{lstlisting}
String lSystem(String axiom, Rule *rules, int ruleCount, int iterations) {
    String current = axiom;
    for (int it = 0; it < iterations; ++it) {
        String next = emptyString();
        for (int i = 0; i < current.length; ++i) {
            bool replaced = false;
            for (int r = 0; r < ruleCount; ++r) {
                if (rules[r].matches(current, i)) {
                    appendString(&next, rules[r].replacement);
                    replaced = true;
                    break;
                }
            }
            if (!replaced) appendChar(&next, current.data[i]);
        }
        current = next;
    }
    return current;
}
\end{lstlisting}

\Needspace{5\baselineskip}
\subsection*{75. Perlin-Style Octave Noise}
\addcontentsline{toc}{subsection}{75. Perlin-Style Octave Noise}

\begin{lstlisting}
float octaveNoise(Vec3 p, int octaves, float persistence, float lacunarity) {
    float total = 0.0f;
    float amplitude = 1.0f;
    float frequency = 1.0f;
    float norm = 0.0f;
    for (int i = 0; i < octaves; ++i) {
        total += amplitude * gradientNoise(scale(p, frequency));
        norm += amplitude;
        amplitude *= persistence;
        frequency *= lacunarity;
    }
    return total / norm;
}
\end{lstlisting}

\Needspace{5\baselineskip}
\subsection*{76. Fault Modifier for Terrain}
\addcontentsline{toc}{subsection}{76. Fault Modifier for Terrain}

\begin{lstlisting}
void applyFault(float **height, int width, int heightN, float displacement) {
    Vec3 p = randomPoint2D(width, heightN);
    Vec3 normal = randomUnit2D();
    for (int y = 0; y < heightN; ++y) {
        for (int x = 0; x < width; ++x) {
            Vec3 q = {x, 0, y};
            float side = dot(sub(q, p), normal);
            height[y][x] += (side > 0.0f) ? displacement : -displacement;
        }
    }
}
\end{lstlisting}

\Needspace{5\baselineskip}
\subsection*{77. Thermal Erosion}
\addcontentsline{toc}{subsection}{77. Thermal Erosion}

\begin{lstlisting}
void thermalErosion(float **h, int w, int n, float talus, float fraction, int iterations) {
    for (int it = 0; it < iterations; ++it) {
        float **delta = zeroGrid(w, n);
        for (int y = 1; y < n-1; ++y) for (int x = 1; x < w-1; ++x) {
            forEachNeighbor4(x, y, nx, ny) {
                float diff = h[y][x] - h[ny][nx];
                if (diff > talus) {
                    float amount = fraction * (diff - talus);
                    delta[y][x] -= amount;
                    delta[ny][nx] += amount;
                }
            }
        }
        addGrid(h, delta, w, n);
    }
}
\end{lstlisting}

\Needspace{5\baselineskip}
\subsection*{78. Hydraulic Erosion}
\addcontentsline{toc}{subsection}{78. Hydraulic Erosion}

\begin{lstlisting}
void hydraulicErosion(Terrain *t, int iterations, float rain, float solubility, float evaporation) {
    for (int it = 0; it < iterations; ++it) {
        addRain(t, rain);
        computeWaterFlow(t);
        forEachCell(t, x, y) {
            float capacity = t->water[y][x] * t->flowSpeed[y][x] * solubility;
            if (t->sediment[y][x] > capacity) deposit(t, x, y, t->sediment[y][x] - capacity);
            else erode(t, x, y, capacity - t->sediment[y][x]);
        }
        moveWaterAndSediment(t);
        evaporateWater(t, evaporation);
    }
}
\end{lstlisting}

\Needspace{5\baselineskip}
\subsection*{79. Maze by Depth-First Backtracking}
\addcontentsline{toc}{subsection}{79. Maze by Depth-First Backtracking}

\begin{lstlisting}
void carveMaze(GridMaze *maze, Cell start) {
    Stack stack = stackCreate();
    markVisited(maze, start);
    stackPush(&stack, start);

    while (!stackEmpty(&stack)) {
        Cell c = stackPeek(&stack);
        Cell choices[4];
        int n = unvisitedNeighbors(maze, c, choices);
        if (n == 0) { stackPop(&stack); continue; }
        Cell next = choices[randomInt(0, n-1)];
        removeWallBetween(maze, c, next);
        markVisited(maze, next);
        stackPush(&stack, next);
    }
}
\end{lstlisting}

\Needspace{5\baselineskip}
\subsection*{80. Dungeon Generate-and-Test Skeleton}
\addcontentsline{toc}{subsection}{80. Dungeon Generate-and-Test Skeleton}

\begin{lstlisting}
Dungeon generateAndTestDungeon(int maxAttempts, bool (*valid)(Dungeon *)) {
    for (int i = 0; i < maxAttempts; ++i) {
        Dungeon d = randomDungeonCandidate();
        if (valid(&d)) return d;
    }
    return repairDungeon(randomDungeonCandidate());
}
\end{lstlisting}

\Needspace{5\baselineskip}
\subsection*{81. Prim Minimum Spanning Tree}
\addcontentsline{toc}{subsection}{81. Prim Minimum Spanning Tree}

\begin{lstlisting}
EdgeList primMST(Graph *g, Node *start) {
    Set visited = setCreate();
    PriorityQueue edges = pqCreate();
    EdgeList tree = edgeListCreate();

    setAdd(&visited, start);
    pushOutgoingEdges(g, start, &edges);

    while (!pqEmpty(&edges) && setSize(&visited) < graphNodeCount(g)) {
        Edge e = pqPopMin(&edges);
        if (setContains(&visited, e.to)) continue;
        edgeListAdd(&tree, e);
        setAdd(&visited, e.to);
        pushOutgoingEdges(g, e.to, &edges);
    }
    return tree;
}
\end{lstlisting}

\Needspace{5\baselineskip}
\subsection*{82. Recursive Subdivision Dungeon/Maze}
\addcontentsline{toc}{subsection}{82. Recursive Subdivision Dungeon/Maze}

\begin{lstlisting}
void subdivide(Space space, int minSize) {
    if (!canSplit(space, minSize)) return;
    Split s = chooseSplit(space);
    Space a, b;
    splitSpace(space, s, &a, &b);
    Corridor c = connectSpaces(a, b);
    addCorridor(c);
    subdivide(a, minSize);
    subdivide(b, minSize);
}
\end{lstlisting}

\Needspace{5\baselineskip}
\subsection*{83. Shape Grammar Execution}
\addcontentsline{toc}{subsection}{83. Shape Grammar Execution}

\begin{lstlisting}
void runShapeGrammar(Grammar *grammar, ObjectList *objects, int maxSteps) {
    for (int step = 0; step < maxSteps; ++step) {
        RuleApplication apps[MAX_RULES];
        int n = triggeredRules(grammar, objects, apps);
        if (n == 0) break;
        RuleApplication chosen = chooseApplication(apps, n);
        chosen.rule->fire(chosen.rule, objects, chosen.target);
    }
}
\end{lstlisting}

\Needspace{5\baselineskip}
\subsection*{84. Shape Grammar Planning}
\addcontentsline{toc}{subsection}{84. Shape Grammar Planning}

\begin{lstlisting}
Plan planGrammar(Grammar *grammar, ObjectList start, bool (*complete)(ObjectList), int maxDepth) {
    Queue q = queueCreate();
    queuePush(&q, (GrammarState){ start, emptyPlan(), 0 });
    while (!queueEmpty(&q)) {
        GrammarState s = queuePop(&q);
        if (complete(s.objects)) return s.plan;
        if (s.depth >= maxDepth) continue;
        RuleApplication apps[MAX_RULES];
        int n = triggeredRules(grammar, &s.objects, apps);
        for (int i = 0; i < n; ++i) {
            ObjectList next = cloneObjects(s.objects);
            apps[i].rule->fire(apps[i].rule, &next, apps[i].target);
            queuePush(&q, (GrammarState){ next, appendPlan(s.plan, apps[i]), s.depth + 1 });
        }
    }
    return emptyPlan();
}
\end{lstlisting}

\medskip\hrule\medskip

\clearpage
\section*{Part VII -- Board-Game Algorithms}
\addcontentsline{toc}{section}{Part VII -- Board-Game Algorithms}

\Needspace{5\baselineskip}
\subsection*{85. Minimax}
\addcontentsline{toc}{subsection}{85. Minimax}

\begin{lstlisting}
float minimax(Board b, int depth, bool maximizing, Move *bestMove) {
    if (depth == 0 || boardGameOver(b)) return boardEvaluate(b);

    Move moves[MAX_MOVES];
    int n = boardMoves(b, moves);
    float best = maximizing ? -INFINITY : INFINITY;

    for (int i = 0; i < n; ++i) {
        Board child = boardMakeMove(b, moves[i]);
        float score = minimax(child, depth - 1, !maximizing, NULL);
        if ((maximizing && score > best) || (!maximizing && score < best)) {
            best = score;
            if (bestMove) *bestMove = moves[i];
        }
    }
    return best;
}
\end{lstlisting}

\Needspace{5\baselineskip}
\subsection*{86. Negamax}
\addcontentsline{toc}{subsection}{86. Negamax}

\begin{lstlisting}
float negamax(Board b, int depth, Move *bestMove) {
    if (depth == 0 || boardGameOver(b)) return boardEvaluateCurrentPlayer(b);

    Move moves[MAX_MOVES];
    int n = boardMoves(b, moves);
    float best = -INFINITY;
    for (int i = 0; i < n; ++i) {
        Board child = boardMakeMove(b, moves[i]);
        float score = -negamax(child, depth - 1, NULL);
        if (score > best) {
            best = score;
            if (bestMove) *bestMove = moves[i];
        }
    }
    return best;
}
\end{lstlisting}

\Needspace{5\baselineskip}
\subsection*{87. Alpha-Beta Negamax}
\addcontentsline{toc}{subsection}{87. Alpha-Beta Negamax}

\begin{lstlisting}
float alphaBetaNegamax(Board b, int depth, float alpha, float beta, Move *bestMove) {
    if (depth == 0 || boardGameOver(b)) return boardEvaluateCurrentPlayer(b);

    Move moves[MAX_MOVES];
    int n = orderedBoardMoves(b, moves);
    for (int i = 0; i < n; ++i) {
        Board child = boardMakeMove(b, moves[i]);
        float score = -alphaBetaNegamax(child, depth - 1, -beta, -alpha, NULL);
        if (score > alpha) {
            alpha = score;
            if (bestMove) *bestMove = moves[i];
        }
        if (alpha >= beta) break;
    }
    return alpha;
}
\end{lstlisting}

\Needspace{5\baselineskip}
\subsection*{88. Aspiration Search}
\addcontentsline{toc}{subsection}{88. Aspiration Search}

\begin{lstlisting}
float aspirationSearch(Board b, int depth, float guess, float window, Move *best) {
    while (true) {
        float alpha = guess - window;
        float beta = guess + window;
        float score = alphaBetaNegamax(b, depth, alpha, beta, best);
        if (score <= alpha) { guess = score; window *= 2.0f; continue; }
        if (score >= beta)  { guess = score; window *= 2.0f; continue; }
        return score;
    }
}
\end{lstlisting}

\Needspace{5\baselineskip}
\subsection*{89. Negascout / Principal Variation Search}
\addcontentsline{toc}{subsection}{89. Negascout / Principal Variation Search}

\begin{lstlisting}
float negascout(Board b, int depth, float alpha, float beta) {
    if (depth == 0 || boardGameOver(b)) return boardEvaluateCurrentPlayer(b);
    Move moves[MAX_MOVES];
    int n = orderedBoardMoves(b, moves);
    float bWindow = beta;

    for (int i = 0; i < n; ++i) {
        Board child = boardMakeMove(b, moves[i]);
        float score = -negascout(child, depth - 1, -bWindow, -alpha);
        if (score > alpha && score < beta && i > 0)
            score = -negascout(child, depth - 1, -beta, -score);
        alpha = fmaxf(alpha, score);
        if (alpha >= beta) break;
        bWindow = alpha + 1.0f;
    }
    return alpha;
}
\end{lstlisting}

\Needspace{5\baselineskip}
\subsection*{90. Zobrist Hashing}
\addcontentsline{toc}{subsection}{90. Zobrist Hashing}

\begin{lstlisting}
uint64_t zobristHash(Board b, uint64_t table[MAX_SQUARES][MAX_PIECES]) {
    uint64_t h = 0;
    for (int sq = 0; sq < MAX_SQUARES; ++sq) {
        int piece = boardPieceAt(b, sq);
        if (piece != EMPTY) h ^= table[sq][piece];
    }
    if (boardSideToMove(b) == BLACK) h ^= SIDE_TO_MOVE_KEY;
    return h;
}
\end{lstlisting}

\Needspace{5\baselineskip}
\subsection*{91. Transposition Table Lookup and Store}
\addcontentsline{toc}{subsection}{91. Transposition Table Lookup and Store}

\begin{lstlisting}
bool ttLookup(TranspositionTable *tt, uint64_t hash, int depth, float alpha, float beta, float *scoreOut) {
    TTEntry *e = tableFind(tt, hash);
    if (!e || e->depth < depth) return false;
    if (e->kind == EXACT) { *scoreOut = e->score; return true; }
    if (e->kind == LOWER_BOUND && e->score >= beta) { *scoreOut = e->score; return true; }
    if (e->kind == UPPER_BOUND && e->score <= alpha) { *scoreOut = e->score; return true; }
    return false;
}

void ttStore(TranspositionTable *tt, uint64_t hash, int depth, float score, BoundKind kind, Move best) {
    TTEntry e = { hash, depth, score, kind, best };
    tableReplace(tt, e);
}
\end{lstlisting}

\Needspace{5\baselineskip}
\subsection*{92. MTD(f)}
\addcontentsline{toc}{subsection}{92. MTD(f)}

\begin{lstlisting}
float mtdf(Board b, int depth, float firstGuess) {
    float g = firstGuess;
    float upper = INFINITY;
    float lower = -INFINITY;
    while (lower < upper) {
        float beta = (g == lower) ? g + 1.0f : g;
        g = alphaBetaNegamax(b, depth, beta - 1.0f, beta, NULL);
        if (g < beta) upper = g;
        else lower = g;
    }
    return g;
}
\end{lstlisting}

\Needspace{5\baselineskip}
\subsection*{93. Monte Carlo Tree Search}
\addcontentsline{toc}{subsection}{93. Monte Carlo Tree Search}

\begin{lstlisting}
typedef struct MCTSNode {
    Board board;
    struct MCTSNode *parent;
    struct MCTSNode **children;
    int childCount;
    int visits;
    float wins;
    Move move;
} MCTSNode;

Move mcts(Board rootBoard, int iterations, float exploration) {
    MCTSNode *root = newMCTSNode(rootBoard, NULL, nullMove());
    for (int i = 0; i < iterations; ++i) {
        MCTSNode *leaf = treePolicy(root, exploration);
        float result = defaultPolicy(leaf->board);
        backup(leaf, result);
    }
    return bestVisitedChild(root)->move;
}

MCTSNode *treePolicy(MCTSNode *node, float c) {
    while (!boardGameOver(node->board)) {
        if (!fullyExpanded(node)) return expand(node);
        node = bestUCTChild(node, c);
    }
    return node;
}

float uctScore(MCTSNode *child, float c) {
    if (child->visits == 0) return INFINITY;
    float exploit = child->wins / child->visits;
    float explore = c * sqrtf(logf(child->parent->visits) / child->visits);
    return exploit + explore;
}
\end{lstlisting}

\Needspace{5\baselineskip}
\subsection*{94. Iterative Deepening Search}
\addcontentsline{toc}{subsection}{94. Iterative Deepening Search}

\begin{lstlisting}
Move iterativeDeepening(Board b, float timeLimitSeconds) {
    Move best = nullMove();
    float guess = 0.0f;
    double start = nowSeconds();
    for (int depth = 1; nowSeconds() - start < timeLimitSeconds; ++depth) {
        Move candidate;
        guess = alphaBetaNegamax(b, depth, -INFINITY, INFINITY, &candidate);
        if (nowSeconds() - start < timeLimitSeconds) best = candidate;
    }
    return best;
}
\end{lstlisting}

\medskip\hrule\medskip

\clearpage
\section*{Part VIII -- Execution Management and World Interfacing}
\addcontentsline{toc}{section}{Part VIII -- Execution Management and World Interfacing}

\Needspace{5\baselineskip}
\subsection*{95. Frequency Scheduler}
\addcontentsline{toc}{subsection}{95. Frequency Scheduler}

\begin{lstlisting}
void runFrequencyScheduler(FrequencyScheduler *sched, GameState *s, float dt) {
    sched->accumulator += dt;
    while (sched->accumulator >= sched->period) {
        sched->accumulator -= sched->period;
        BehaviorRecord *r = &sched->records[sched->cursor];
        r->behavior->run(r->behavior, s, sched->period);
        sched->cursor = (sched->cursor + 1) % sched->recordCount;
    }
}
\end{lstlisting}

\Needspace{5\baselineskip}
\subsection*{96. Load-Balancing Scheduler}
\addcontentsline{toc}{subsection}{96. Load-Balancing Scheduler}

\begin{lstlisting}
void runLoadBalancingScheduler(LoadScheduler *sched, GameState *s, float budgetMs) {
    double start = nowMilliseconds();
    while (nowMilliseconds() - start < budgetMs && sched->hasWork(sched)) {
        BehaviorRecord *r = schedNextBehavior(sched);
        r->behavior->run(r->behavior, s, r->sliceSeconds);
        updateEstimatedCost(sched, r, nowMilliseconds() - start);
    }
}
\end{lstlisting}

\Needspace{5\baselineskip}
\subsection*{97. Priority Scheduler}
\addcontentsline{toc}{subsection}{97. Priority Scheduler}

\begin{lstlisting}
void runPriorityScheduler(PriorityScheduler *sched, GameState *s, float budgetMs) {
    double start = nowMilliseconds();
    rebuildPriorityQueue(sched);
    while (!pqEmpty(&sched->queue) && nowMilliseconds() - start < budgetMs) {
        BehaviorRecord *r = pqPopMax(&sched->queue);
        r->behavior->run(r->behavior, s, r->allocatedTime);
        r->lastRunTime = nowSeconds();
    }
}
\end{lstlisting}

\Needspace{5\baselineskip}
\subsection*{98. Behavioral Level of Detail}
\addcontentsline{toc}{subsection}{98. Behavioral Level of Detail}

\begin{lstlisting}
void runBehaviorLOD(BehaviorLOD *lod, GameState *s, Character *c, float dt) {
    float importance = computeImportance(c, s);
    for (int i = 0; i < lod->levelCount; ++i) {
        if (importance >= lod->levels[i].enterAt && importance <= lod->levels[i].exitBelow) {
            if (lod->current != i) {
                if (lod->current >= 0) lod->levels[lod->current].exit(c, s);
                lod->current = i;
                lod->levels[i].enter(c, s);
            }
            lod->levels[i].behavior->run(lod->levels[i].behavior, s, dt);
            return;
        }
    }
}
\end{lstlisting}

\Needspace{5\baselineskip}
\subsection*{99. Event Manager}
\addcontentsline{toc}{subsection}{99. Event Manager}

\begin{lstlisting}
void dispatchEvents(EventManager *mgr, GameState *s) {
    while (!queueEmpty(&mgr->events)) {
        Event e = queuePop(&mgr->events);
        for (int i = 0; i < mgr->listenerCount; ++i) {
            ListenerRegistration *r = &mgr->listeners[i];
            if (r->type == e.type && r->filter(&e, r->ctx)) r->listener->notify(r->listener, &e, s);
        }
    }
}
\end{lstlisting}

\Needspace{5\baselineskip}
\subsection*{100. Polling Station with Cached Tasks}
\addcontentsline{toc}{subsection}{100. Polling Station with Cached Tasks}

\begin{lstlisting}
void pollStation(PollingStation *station, GameState *s, float now) {
    for (int i = 0; i < station->taskCount; ++i) {
        PollingTask *t = &station->tasks[i];
        if (now - t->lastUpdate >= t->refreshInterval) {
            t->cachedValue = t->compute(s, t->ctx);
            t->lastUpdate = now;
        }
    }
}

Value pollingGet(PollingStation *station, int taskId, GameState *s, float now) {
    PollingTask *t = &station->tasks[taskId];
    if (now - t->lastUpdate >= t->refreshInterval) {
        t->cachedValue = t->compute(s, t->ctx);
        t->lastUpdate = now;
    }
    return t->cachedValue;
}
\end{lstlisting}

\Needspace{5\baselineskip}
\subsection*{101. Regional Sense Manager}
\addcontentsline{toc}{subsection}{101. Regional Sense Manager}

\begin{lstlisting}
void addSenseSignal(RegionalSenseManager *mgr, Signal sig) {
    Region *r = regionForPosition(mgr, sig.position);
    regionSignalListAdd(r, sig);
}

void sendSenseSignals(RegionalSenseManager *mgr, Sensor *sensors, int sensorCount) {
    for (int i = 0; i < sensorCount; ++i) {
        RegionSet nearby = regionsNear(mgr, sensors[i].position, sensors[i].range);
        forEachRegion(nearby, r) {
            forEachSignal(r->signals, sig) {
                if (sensors[i].detects(&sensors[i], sig.modality) && modalityExtraChecks(sig, sensors[i])) {
                    sensors[i].notify(&sensors[i], sig);
                }
            }
        }
    }
    clearExpiredSignals(mgr);
}
\end{lstlisting}

\medskip\hrule\medskip

\clearpage
\section*{Implementation Notes}
\addcontentsline{toc}{section}{Implementation Notes}

\begin{enumerate}
\item \textbf{Game AI favors controllability over theoretical purity.} Many algorithms above are intentionally approximate; their job is to produce believable behavior under frame-time constraints.
\item \textbf{Data representation matters as much as the algorithm.} Pathfinding, perception, and tactical analysis depend on world representations such as tile graphs, navmeshes, waypoint graphs, influence grids, and region partitions.
\item \textbf{Scheduling is part of the AI algorithm.} Expensive algorithms should be interruptible, cached, distributed across frames, or degraded by level of detail.
\item \textbf{Composition is the normal production pattern.} Steering outputs are blended or prioritized; decisions are layered through FSMs, behavior trees, utility selectors, blackboards, and planners; pathfinding is coupled to tactical costs and movement controllers.
\item \textbf{Learning is usually offline or constrained.} The learning algorithms are valuable, but production games often use them selectively because online learning can reduce designer control and introduce unstable behavior.
\end{enumerate}

\medskip\hrule\medskip

\clearpage
\section*{Algorithm Index}
\addcontentsline{toc}{section}{Algorithm Index}

\begin{longtable}{L{0.24\textwidth}L{0.68\textwidth}}
\toprule
Area & Algorithms \\
\midrule\endfirsthead
\toprule
Area & Algorithms \\
\midrule\endhead
Movement & Kinematic seek/flee/arrive/wander; dynamic seek/flee/arrive/align/velocity match; pursue/evade; face/look-where-going; wander; path following; separation; collision avoidance; obstacle avoidance; blended steering; priority steering; steering pipeline; firing solution; iterative targeting; jumping; formation slots; 3D orientation variants. \\
Pathfinding & Dijkstra; A\textsuperscript{*}; node-array A\textsuperscript{*}; path smoothing; hierarchical pathfinding; instanced-geometry reuse; continuous-time pathfinding; movement planning. \\
Decision Making & Decision trees; FSM; HSM; behavior trees; decorators; parallel behavior trees; fuzzy logic; fuzzy FSM; Markov FSM; goal-oriented behavior; GOAP; rule systems; Rete; blackboard; action manager. \\
Tactical AI & Cover scoring; waypoint condensation; influence maps; map flooding; convolution; separable filters; cellular automata; tactical path-cost blending. \\
Learning & Hill climbing; simulated annealing; N-grams; hierarchical N-grams; Naive Bayes; ID3; continuous splits; Q-learning; MLP feedforward/backpropagation. \\
PCG & Seed mixing; Halton sequence; phyllotaxis; Poisson disk; L-systems; octave noise; terrain faults; thermal erosion; hydraulic erosion; DFS maze; generate-and-test; Prim MST; recursive subdivision; shape grammars; grammar planning. \\
Board Games & Minimax; negamax; alpha-beta; aspiration windows; negascout/PVS; Zobrist hashing; transposition tables; MTD(f); MCTS; iterative deepening. \\
Infrastructure & Frequency scheduler; load-balancing scheduler; priority scheduler; behavioral LOD; event manager; polling station; regional sense manager. \\
\bottomrule
\end{longtable}

\clearpage
\section*{Document Control Notes}
\phantomsection
\addcontentsline{toc}{section}{Document Control Notes}
\begin{itemize}
\item The algorithms are independent C-like reconstructions intended for study, architecture review, engine design, and implementation planning.
\item The notation normalizes game-AI vocabulary across movement, steering, pathfinding, decision making, tactical analysis, learning, procedural generation, game-tree search, scheduling, and world-interface systems.
\item For production use, each pseudocode routine should be paired with deterministic tests, frame-budget measurements, data-representation validation, debug visualizations, and gameplay-tuning hooks.
\end{itemize}
\end{document}
